\documentclass[12pt]{article}
\usepackage[utf8]{inputenc}
\usepackage[T1]{fontenc}
\usepackage{hyperref}
\title{Principense Language Resources for K-12 Curriculum Development}
\author{Compiled by Claude Sonnet 5.5 (claude-sonnet-5-5) for Andr\'as Kornai}
\date{}
\begin{document}
\maketitle

\section{Introduction}
Principense (Lung'Ie, also lungw'Ie or lingw'Ie) is a Portuguese-based creole of the island of Pr\'incipe, S\~ao Tom\'e and Pr\'incipe. It is endangered. Estimates of native speakers range from under one hundred to about 200 (Ethnologue, 2021).

\section{Primary Educational Resources}
\begin{itemize}
\item Lung'Ie online dictionary (Living Dictionaries), from research led by Ana L\'ivia Agostinho (UFSC), with speaker recordings; its stated purpose is supporting Lung'Ie teaching in Pr\'incipe schools: \url{https://livingdictionaries.app/lungie/about}
\end{itemize}
Agostinho's doctoral work also contains a pedagogical grammar of seventeen lessons with a Lung'Ie/Portuguese glossary of about 1,200 words. Search results did not give a direct link to it.

\section{Linguistic and Reference Materials}
\begin{itemize}
\item APiCS survey chapter on Principense: \url{https://apics-online.info/surveys/37}
\item Wikipedia, Principense Creole: \url{https://en.wikipedia.org/wiki/Principense_Creole}
\item LINGUIST List review (issue 28.2786): \url{https://linguistlist.org/issues/28/2786/}
\item Journal article in Finnish Journal of Linguistics: \url{https://journal.fi/finjol/article/view/152996/98520}
\end{itemize}
Print works named in search results, without links: Philippe Maurer, \emph{Principense (Lung'Ie): Grammar, Texts, and Vocabulary} (2009; includes an English word list of 1650 entries and an edition of Ribeiro 1888); Agostinho's phonology (LINCOM, 2016); G\"unther (1973); Pinheiro de Araujo (2012) Principense--Portuguese dictionary (master's thesis).

\section{Religious and Community Resources}
None identified in the search results.

\section{Digital and Audio Resources}
The Living Dictionaries site above includes audio recordings by speakers.

\section{Curriculum Development Context}
Documentation is good for a language of this size (a full grammar, a dictionary, a pedagogical grammar), but no school curriculum, primers or readers were found. With few elderly speakers, recording work should be urgent.

\end{document}